\section{Exponential positivity for irreducible CP maps}


\begin{theorem}[Exponential positivity for irreducible CP maps]
    \label{thm:exp_pd_irred}
    \lean{exp_posDef_of_irreducible_cp}
    \leanok
    \uses{def:irreducible_cp}
    Let $E$ be an irreducible completely positive map on $\MN{D}$,
    let $A \ge 0$ be nonzero, and let $t > 0$.
    Then
    \begin{align}
        \exp(tE)(A)
         & = \sum_{k=0}^{\infty}\frac{t^k}{k!}E^k(A)>0.
        \label{eq:qpf_exp_positive}
    \end{align}
    This is the completely positive specialization of the forward implication
    in~\cite[Theorem~6.2(3)]{Wolf2012Quantum}.
\end{theorem}

\begin{proof}
    \leanok
    \uses{thm:exp_truncation_pd_irred}
    By Theorem~\ref{thm:exp_truncation_pd_irred}, the first $D$ terms in
    (\ref{eq:qpf_exp_positive}) already form a positive definite matrix.
    Every remaining term is positive semidefinite, so the full exponential
    series is the sum of a positive definite matrix and a positive
    semidefinite tail.
\end{proof}

\begin{theorem}[Exponential characterization of irreducibility]
    \label{thm:wolf_6_2_item_3}
    \lean{irreducible_iff_exp_posDef_forall}
    \leanok
    \uses{def:irreducible_cp}
    Let $E$ be a completely positive map on $\MN{D}$. Then $E$ is
    irreducible if and only if, for every $t>0$ and every nonzero
    $A\geq0$, one has $\exp(tE)(A)>0$.
    This is the completely positive specialization of the equivalence in
    \cite[Theorem~6.2(3)]{Wolf2012Quantum};
    Theorem~\ref{thm:exp_pd_irred} is its forward implication.
\end{theorem}

\begin{proof}\leanok
    \uses{thm:exp_pd_irred}
    The forward implication is Theorem~\ref{thm:exp_pd_irred}. Conversely,
    let $P$ be an invariant orthogonal projection. Invariance gives
    $E(P\MN{D}P)\subseteq P\MN{D}P$, so induction yields
    $E^k(P)\in P\MN{D}P$ for every $k\geq0$. Therefore every term, and hence
    the convergent series
    \begin{align}
        \exp(tE)(P)
         & =\sum_{k=0}^{\infty}\frac{t^k}{k!}E^k(P),
        \notag
    \end{align}
    lies in the same corner. If $P\neq0$, then at $t=1$ the assumed positive
    definiteness is impossible unless $P=\Id$, because a positive definite
    matrix cannot be supported on a proper corner. Thus every invariant
    projection is $0$ or $\Id$, so $E$ is irreducible.
\end{proof}

\section{Ergodicity of irreducible channels}

\begin{theorem}[Unique density-matrix fixed point for an irreducible channel]
    \label{thm:channel_density_fp_irred}
    \lean{IsChannel.exists_unique_density_fixedPoint_of_irreducible}
    \leanok
    \uses{def:channel, def:irreducible_cp}
    Let $E$ be an irreducible quantum channel on $\MN{D}$ with $D > 0$.
    Then there exists a unique density matrix $\sigma$ such that
    $E(\sigma) = \sigma$.
    The fixed point $\sigma$ is positive definite.
    This is the fixed-point conclusion in the forward direction of
    \cite[Corollary~6.3]{Wolf2012Quantum}, specialized to completely
    positive trace-preserving maps.
\end{theorem}

\begin{proof}
    \leanok
    \uses{lem:cesaro_mean_density, lem:cesaro_subseq_fixed,
        thm:psd_eigenvector_pd_irred, thm:psd_eigenvector_unique_irred}
    Every Ces\`aro mean (\ref{eq:channel_cesaro_mean}) of a density matrix
    stays inside the compact convex set of density matrices, so a subsequence
    converges to a density matrix~$\sigma$.  The telescope identity
    (\ref{eq:channel_cesaro_telescope}) gives
    $E(\sigma)=\sigma$.  Theorem~\ref{thm:psd_eigenvector_pd_irred}
    makes $\sigma$ positive definite.  If $\tau$ is another density-matrix
    fixed point, Theorem~\ref{thm:psd_eigenvector_unique_irred} gives
    $\tau=c\sigma$.  Taking traces yields $1=c$, hence $\tau=\sigma$.
\end{proof}

\begin{theorem}[Ces\`aro convergence for irreducible channels]
    \label{thm:channel_cesaro_irred}
    \lean{IsChannel.cesaroMean_tendsto_of_irreducible}
    \leanok
    \uses{def:channel, def:irreducible_cp, def:cesaro_mean}
    Let $E$ be an irreducible quantum channel on $\MN{D}$ with $D > 0$,
    and let $\rho$ be any density matrix.
    Then the Ces\`aro means (\ref{eq:channel_cesaro_mean}) converge to the
    unique positive definite density-matrix fixed point of~$E$.
    This is the Ces\`aro-convergence conclusion in the forward direction of
    \cite[Corollary~6.3]{Wolf2012Quantum}, specialized to quantum channels.
\end{theorem}

\begin{proof}
    \leanok
    \uses{thm:channel_density_fp_irred}
    Every subsequential limit of the Ces\`aro means is again a density-matrix
    fixed point.
    By Theorem~\ref{thm:channel_density_fp_irred}, there is only one such
    fixed point, so every convergent subsequence has the same limit.
    Compactness then forces the whole sequence to converge to that limit.
\end{proof}

\section{A CP spectral characterization of irreducibility}

\begin{definition}[Restricted CP spectral properties]
    \label{def:has_spectral_properties}
    \lean{HasSpectralProperties}
    \leanok
    \uses{def:cp_map}
    A completely positive map $E$ on $\MN{D}$ has the \emph{restricted CP
        spectral properties} used here if there exist a positive real number $r$, a positive
    definite right eigenvector $\rho$, and a positive definite left eigenvector
    $\sigma$ for the adjoint map such that
    \begin{align}
        E(\rho)
         & =r\rho, \notag \\
        E^\dagger(\sigma)
         & =r\sigma,
        \label{eq:qpf_spectral_eigenvectors}
    \end{align}
    every positive semidefinite right eigenvector for eigenvalue $r$ is a
    scalar multiple of $\rho$, and the spectral radius of $E$ is equal to~$r$.
    The uniqueness clause concerns only positive semidefinite Perron
    eigenvectors. Unlike the full nondegenerate-eigenspace statement in
    \cite[Theorem~6.4]{Wolf2012Quantum}, it does not assert that every complex
    eigenvector at $r$ is proportional to $\rho$.
\end{definition}

\begin{theorem}[Irreducibility gives the restricted CP spectral properties]
    \label{thm:has_spectral_properties_irred}
    \lean{hasSpectralProperties_of_irreducible_cp}
    \leanok
    \uses{def:has_spectral_properties, def:irreducible_cp}
    Every nonzero irreducible completely positive map on $\MN{D}$ has the
    restricted CP spectral properties of Definition~\ref{def:has_spectral_properties}.
\end{theorem}

\begin{proof}
    \leanok
    \uses{thm:pf_eigenvector_existence, thm:psd_eigenvector_unique_irred,
        thm:spectral_radius_pf_irred}
    Combine Perron--Frobenius existence for the map and its adjoint with the
    uniqueness theorem for positive semidefinite Perron eigenvectors and the
    spectral-radius identity of
    Theorem~\ref{thm:spectral_radius_pf_irred}.
\end{proof}

\begin{theorem}[Restricted CP spectral properties imply irreducibility]
    \label{thm:irred_from_spectral_properties}
    \lean{isIrreducibleMap_of_hasSpectralProperties}
    \leanok
    \uses{def:has_spectral_properties}
    If a completely positive map on $\MN{D}$ has the restricted CP spectral properties of
    Definition~\ref{def:has_spectral_properties}, then it is irreducible.
\end{theorem}

\begin{proof}
    \leanok
    \uses{thm:tp_gauge_from_adjoint_fixed, lem:cesaro_subseq_fixed}
    Gauge by the positive definite left eigenvector in
    (\ref{eq:qpf_spectral_eigenvectors}) and rescale by the Perron eigenvalue
    to obtain a trace-preserving map with a positive definite fixed point.
    A nontrivial invariant projection would then produce, by Ces\`aro averaging
    in its corner algebra, a second positive semidefinite fixed point not
    proportional to the Perron vector, contradicting the uniqueness clause in
    Definition~\ref{def:has_spectral_properties}.
\end{proof}

\begin{theorem}[Irreducibility iff the restricted CP spectral properties hold]
    \label{thm:irred_iff_spectral_properties}
    \lean{isIrreducibleMap_iff_spectral_properties}
    \leanok
    \uses{def:has_spectral_properties, def:irreducible_cp}
    Let $E$ be a nonzero completely positive map on $\MN{D}$.
    Then $E$ is irreducible if and only if it has the restricted CP spectral properties of
    Definition~\ref{def:has_spectral_properties}. This is a CP-map variant of
    \cite[Theorem~6.4]{Wolf2012Quantum}: uniqueness is required only among
    positive semidefinite Perron eigenvectors, not on the full eigenspace.
    Those restricted properties are nevertheless sufficient for the
    irreducibility equivalence.
\end{theorem}

\begin{proof}
    \leanok
    \uses{thm:has_spectral_properties_irred, thm:irred_from_spectral_properties}
    Combine Theorems~\ref{thm:has_spectral_properties_irred}
    and~\ref{thm:irred_from_spectral_properties}.
\end{proof}
